%%%%%%%%%%  scan idx 489  |  folio 476  %%%%%%%%%%
\SGAPrintPage{476}%
\sgaStableTarget{sga6.e08.prop.2-2}\textbf{Proposition 2.2.} \uline{The complex} $\underline{L}^{X/Y}_\bullet$ \uline{does not depend}, \uline{up to canonical isomorphism in the derived category} $D(X)$, \uline{on the chosen factorization through a smooth} $Y$-\uline{scheme}.

Let $V,V'$ be two smooth $Y$-schemes through which $f$ factors. Put: $V''=V\times_Y V'$; then $V''$ is smooth over $Y$, and also over $V$ and $V'$. Moreover, the immersions $X\to V$ and $X\to V'$ define an immersion $X\to V''$ which factors $f$. One may therefore reduce to comparing the complexes $\underline{L}^{X/Y}_\bullet$ associated to two factorizations through smooth $Y$-schemes $V$ and $V'$, such that there exists a commutative diagram over $Y$:
\[
\begin{tikzcd}[column sep=large,row sep=large]
& V' \arrow[d,"p"]\\
X \arrow[ur,"j"] \arrow[r,"i"] & V,
\end{tikzcd}
\]
where $i$ and $j$ are closed immersions, and $p$ a smooth morphism.

By functoriality, one then obtains homomorphisms:
\[
\underline{I}/\underline{I}^2\longrightarrow \underline{J}/\underline{J}^2;
\qquad
p^*(\underline{\Omega}^1_{V/Y})\longrightarrow\underline{\Omega}^1_{V'/Y},
\]
hence a diagram:
\sgaStableTarget{sga6.e08.eq.2-2-1}\begin{equation}\tag{2.2.1}
\begin{tikzcd}[column sep=large,row sep=large]
\underline{I}/\underline{I}^2 \arrow[r] \arrow[d,"d"'] & \underline{J}/\underline{J}^2 \arrow[d,"d"]\\
i^*(\underline{\Omega}^1_{V/Y}) \arrow[r] & j^*(\underline{\Omega}^1_{V'/Y}).
\end{tikzcd}
\end{equation}

Since $V$ and $V'$ are smooth over $Y$, one obtains the exact sequence:
\[
0\longrightarrow p^*(\underline{\Omega}^1_{V/Y})\longrightarrow\underline{\Omega}^1_{V'/Y}
\longrightarrow\underline{\Omega}^1_{V'/V}\longrightarrow0.
\]
Since the Modules of differentials considered are locally free over $\underline{O}_{V'}$, applying the functor $j^*$ gives the exact sequence:
\sgaStableTarget{sga6.e08.eq.2-2-2}\begin{equation}\tag{2.2.2}
0\longrightarrow i^*(\underline{\Omega}^1_{V/Y})\longrightarrow j^*(\underline{\Omega}^1_{V'/Y})
\longrightarrow j^*(\underline{\Omega}^1_{V'/V})\longrightarrow0.
\end{equation}

%%%%%%%%%%  scan idx 490  |  folio 477  %%%%%%%%%%
\SGAPrintPage{477}%
Furthermore,
$j$ defines a section $j'$ of $X'=V'\times_V X$:
\[
\begin{tikzcd}[column sep=large,row sep=large]
X' \arrow[r,"i'"] \arrow[d,"p'"'] & V' \arrow[d,"p"]\\
X \arrow[r,"i"] \arrow[ur,"j"] \arrow[u,"j'"',bend right] & V .
\end{tikzcd}
\]
Let $\underline{I},\underline{J},\underline{K}$ be the ideals defining $i,j,j'$; we have: $j=i'\circ j'$, and the immersion $j'$ is regular; moreover, $i'$ is defined by the ideal $p^*(\underline{I})$ (since $p$ is flat), and its conormal sheaf is therefore $i'^*(p^*(\underline{I}))=p'^*(\underline{I}/\underline{I}^2)$. One then obtains the exact sequence:
\sgaStableTarget{sga6.e08.eq.2-2-3}\begin{equation}\tag{2.2.3}
0\longrightarrow \underline{I}/\underline{I}^2\longrightarrow \underline{J}/\underline{J}^2\longrightarrow \underline{K}/\underline{K}^2\longrightarrow0.
\end{equation}
We shall show:

a) that the diagram (2.2.1) is commutative; this will thus give a morphism of complexes $\theta$ from the complex $\underline{L}^{X/Y}_\bullet$ defined by $V$ to the complex $\underline{L}^{X/Y}_\bullet$ defined by $V'$;

b) that the homomorphism: $\underline{K}/\underline{K}^2\longrightarrow j^*(\underline{\Omega}^1_{V'/V})$ defined by (2.2.2), (2.2.3) and (2.2.1) is an isomorphism; this will imply that the mapping cylinder of $\theta$ is acyclic, hence that $\theta$ is a quasi-isomorphism, and therefore defines an isomorphism in the derived category.

The properties a) and b) are local on $X$. One may therefore suppose $X,Y,V$ and $V'$ affine. Put: $Y=\Spec(A)$, $V=\Spec(B)$, $V'=\Spec(B')$, $X=\Spec(C)$, and keep the same letters to denote the corresponding homomorphisms and ideals. We have the diagram:

%%%%%%%%%%  scan idx 491  |  folio 478  %%%%%%%%%%
\SGAPrintPage{478}%
\[
\begin{tikzcd}[column sep=large,row sep=large]
B\otimes_A B \arrow[r,"p\otimes p"] \arrow[d,"\Delta"'] \arrow[dr,phantom,"(1)"] & B'\otimes_A B' \arrow[d,"\Delta"]\\
B \arrow[r,"p"] \arrow[dr,"i"'] & B' \arrow[d,"j"'] \arrow[dr,"i'","(2)"' near start] & B'\otimes_B B' \arrow[l,"\Delta"'] \arrow[d]\\
{} & C \arrow[r,shift right=.7ex,"p'"'] & B'\otimes_B C = B'/IB'. \arrow[l,shift right=.7ex,"j'"']
\end{tikzcd}
\]

The homomorphism $\underline{I}/\underline{I}^2\to\underline{J}/\underline{J}^2$ is obtained from the homomorphism: $IB'\to J$ and coincides with the one obtained from $p:I\to J$; the homomorphism: $i^*(\underline{\Omega}^1_{V/Y})\longrightarrow j^*(\underline{\Omega}^1_{V'/Y})$ is obtained from the commutative square (1). Property a) is therefore evident.

The homomorphism: $\underline{K}/\underline{K}^2\to j^*(\underline{\Omega}^1_{V'/V})$ considered in b) can be defined as follows: let $\overline x\in K/K^2$, and let $x$ be an element of $J$ lifting $\overline x$ (one exists since $K=J/IB'$). To $\overline x$ one associates the image of $x$ under the composite homomorphism: $J\xrightarrow{d}\Omega_{B'/A}\otimes_{B'}C\longrightarrow\Omega^1_{B'/B}\otimes_{B'}C$. Furthermore, let $H$ be the augmentation ideal of: $B'\otimes_B B'\longrightarrow B'$; the commutative square (2) gives a homomorphism: $H\to K$, hence a homomorphism: $H/H^2\otimes_{B'}C\longrightarrow K/K^2$. Now $H/H^2=\Omega^1_{B'/B}$; consequently, it suffices to check that the homomorphisms constructed are inverse to each other, which is done without difficulty.

\sgaStableTarget{sga6.e08.def.2-3}\textbf{Definition 2.3.} \uline{Let} $f:X\to Y$ \uline{be a smoothable morphism}. \uline{The complex} $\underline{L}^{X/Y}_\bullet$, \uline{defined up to canonical isomorphism in the derived category} $D(X)$, \uline{is called the} relative cotangent complex \uline{of} $X$ \uline{over} $Y$ (\uline{or of} $f$).

\sgaStableTarget{sga6.e08.prop.2-4}\textbf{Proposition 2.4.} \uline{Let} $f:X\to Y$ \uline{be a complete-intersection morphism} (1.1) \uline{which is smoothable}. \uline{Then} $\underline{L}^{X/Y}_\bullet$ \uline{is perfect}. \uline{More precisely}, \uline{for every}

%%%%%%%%%%  scan idx 492  |  folio 479  %%%%%%%%%%
\SGAPrintPage{479}%
\uline{factorization of} $f$ \uline{as the composite of an immersion and a smooth morphism}, \uline{the complex} $\underline{L}^{X/Y}_\bullet$ \uline{defined in} 2.1 \uline{is strictly perfect}.

Let $f=g\circ i$, where $g$ is a smooth morphism, and $i$ an immersion. By 1.2, $i$ is a regular immersion. If $\underline{J}$ is the Ideal defining $i$, then $\underline{J}/\underline{J}^2$ is a locally free $\underline{O}_X$-Module of finite type. Similarly, since $g:V\to Y$ is smooth, $i^*(\underline{\Omega}^1_{V/Y})$ is locally free of finite type, giving 2.4.

\sgaStableTarget{sga6.e08.cor.2-5}\textbf{Corollary 2.5.} \uline{Let} $f:X\to Y$ \uline{be a smoothable complete-intersection morphism}, \uline{and} $f=g\circ i$ \uline{a factorization of} $f$ \uline{through a smooth} $Y$-\uline{scheme} $V$. \uline{Let} $\underline{J}$ \uline{be the ideal of} $\underline{O}_V$ \uline{defining} $i$. \uline{Then}, \uline{if} $K^{\bullet}(X)$ \uline{denotes the Grothendieck group of locally free} $\underline{O}_X$-\uline{Modules of finite type} (\uline{resp}. \uline{of perfect complexes on} $X$), \uline{the element} $T_f\in K^{\bullet}(X)$ \uline{defined by}:
\[
T_f=\operatorname{cl}(i^*(\underline{\Omega}^1_{V/Y}))-\operatorname{cl}(\underline{J}/\underline{J}^2)
\]
\[
(\text{resp. }T_f=\operatorname{cl}(\underline{L}^{X/Y}_\bullet))
\]
\uline{does not depend on the chosen factorization of} $f$. \uline{The value at a point} $x$ \uline{of} $X$ \uline{of the augmentation of} $T_f$ \uline{is precisely the virtual relative dimension of} $X$ \uline{over} $Y$ \uline{at the point} $x$ (1.9).

The element $T_f$ is well defined in both cases considered by 2.4, and its class in $K^{\bullet}(X)$ is independent of the factorization, by 2.2 in the second case, and by the proof of 2.2 in the first case. The last assertion of 2.5 is trivial.

\sgaStableTarget{sga6.e08.prop.2-6}\textbf{Proposition 2.6.} \uline{Let} $f:X\to Y$, \uline{and} $g:Y\to Z$ \uline{be two smoothable complete-intersection morphisms}, \uline{and suppose in addition that the following condition is satisfied}:

%%%%%%%%%%  scan idx 493  |  folio 480  %%%%%%%%%%
\SGAPrintPage{480}%
(G) \uline{There
exists a smooth} $Y$-\uline{scheme} $V$, \uline{factoring} $f$, \uline{such that for every closed immersion}: $Y\to V'$ \uline{one can find a smooth} $V'$-\uline{scheme} $V''$, \uline{with} $V\simeq V''\times_{V'}Y$ (\uline{this condition is satisfied in particular if} $f$ \uline{is quasi-projective and if} $Y$ \uline{has an ample sheaf}, \uline{by taking for} $V$ \uline{a} $\mathbf{P}^r_Y$).

\uline{Then} $g\circ f$ \uline{is a complete intersection}, \uline{is smoothable}, \uline{and we have a distinguished triangle}:
\[
\begin{tikzcd}[column sep=large,row sep=large]
& \underline{L}^{X/Y}_\bullet \arrow[dl] \\
\mathbb Lf^*(\underline{L}^{Y/Z}_\bullet) \arrow[rr] && \underline{L}^{X/Z}_\bullet . \arrow[ul]
\end{tikzcd}
\]

Let $V'$ be a smooth $Z$-scheme factoring $g$. Consider the commutative diagram:
\[
\begin{tikzcd}[column sep=large,row sep=large]
X \arrow[r,"i"] \arrow[dr,"f"'] & V \arrow[r,"j'"] \arrow[d,"p"] & V'' \arrow[d,"p''"]\\
& Y \arrow[r,"j"] \arrow[dr,"g"'] & V' \arrow[d,"p'"]\\
& & Z .
\end{tikzcd}
\]

Then the immersion $j'$ is regular, since it is obtained from $j$ by a flat base change. Consequently $j'\circ i$ is a regular immersion factoring $g\circ f$, which is therefore smoothable and a complete intersection. Moreover, we have the exact sequence of conormal sheaves:
\[
0\longrightarrow i^*(\underline{J}'/\underline{J}'^2)\longrightarrow \underline{K}/\underline{K}^2\longrightarrow \underline{I}/\underline{I}^2\longrightarrow0,
\]
where $\underline{I},\underline{J},\underline{J}',\underline{K}$ denote the sheaves defining respectively the immersions $i,j,j',j'\circ i$. Now, since $p''$ is flat, $\underline{J}'=p''^*(\underline{J})$. The sequence can therefore be written:
\sgaStableTarget{sga6.e08.eq.2-6-1}\begin{equation}\tag{2.6.1}
0\longrightarrow f^*(\underline{J}/\underline{J}^2)\longrightarrow \underline{K}/\underline{K}^2\longrightarrow \underline{I}/\underline{I}^2\longrightarrow0.
\end{equation}

Furthermore, since $V'$ and $V''$ are smooth over $Z$, we have the exact sequence:
\[
0\longrightarrow p''^*(\underline{\Omega}^1_{V'/Z})\longrightarrow\underline{\Omega}^1_{V''/Z}\longrightarrow\underline{\Omega}^1_{V''/V'}\longrightarrow0.
\]
